\clearpage\pagestyle{otMain}\setlength{\otBlocEspacement}{2.5mm}\fontsize{11}{15}\selectfont\corrigesChapitre{5}{Égalités, inconnues et régularités}
\begin{blocLivre}{cor-5-1}\begin{corrige}{1}{Vrai ou faux ?}\demarche{Calculer séparément les deux côtés.}
$8+6=14$ et $9+5=14$ : la première égalité est $\res{\text{vraie}}$. $7\times4=28$ mais $30-3=27$ : la seconde est $\res{\text{fausse}}$.
\end{corrige}\end{blocLivre}
\begin{blocLivre}{cor-5-2}\begin{corrige}{2}{Trous additifs}\demarche{Utiliser l’opération inverse.}
$92-47=45$, donc $\square=\res{45}$. Pour la seconde : $35+28=63$, donc $\square=\res{63}$.
\end{corrige}\end{blocLivre}
\begin{blocLivre}{cor-5-3}\begin{corrige}{3}{Trous multiplicatifs}\demarche{Commencer par le côté entièrement connu.}
$84:7=12$, donc $\square=\res{12}$. À droite, $6\times8=48$ ; alors $178-\square=48$ donne $\square=178-48=\res{130}$.
\end{corrige}\end{blocLivre}
\begin{blocLivre}{cor-5-4}\begin{corrige}{4}{La machine à nombres}\demarche{Suivre les étapes dans l’ordre.}
Avec 6 : $6+2=8$ puis $8\times4=\res{32}$. Avec 10 : $10+2=12$ puis $12\times4=\res{48}$.
\label{fig-cor-5-4}\figureGeo[corrige]{prog-4}
\end{corrige}\end{blocLivre}
\begin{blocLivre}{cor-5-5}\begin{corrige}{5}{Poursuivre une suite}\demarche{Calculer l’écart entre deux termes voisins.}
Chaque terme vaut le précédent $\res{\text{plus }7}$. Les termes suivants sont $24+7=\res{31}$ puis $31+7=\res{38}$.
\end{corrige}\end{blocLivre}
\begin{blocLivre}{cor-5-6}\begin{corrige}{6}{Deux règles différentes}\demarche{Tester « plus combien » puis « fois combien ».}
Dans la première, chaque terme est $\res{\text{le double}}$ du précédent : le suivant est $\res{80}$. Dans la seconde, on ajoute 5 : le suivant est $\res{25}$. Deux suites peuvent commencer pareil et diverger.
\end{corrige}\end{blocLivre}
\begin{blocLivre}{cor-5-7}\begin{corrige}{7}{Remonter la machine}\demarche{Défaire les étapes en sens inverse.}
On remonte à l’envers : avant de retirer 4, le nombre était $20+4=24$ ; avant de multiplier par 3, il était $24:3=\res8$. Contrôle : $8\times3-4=20$.
\end{corrige}\end{blocLivre}
\begin{blocLivre}{cor-5-8}\begin{corrige}{8}{La suite qui zigzague}\demarche{Observer les écarts un pas sur deux.}
La règle alterne : $+5$ puis $-2$. Après 89 : $89+5=\res{94}$, puis $94-2=\res{92}$.
\end{corrige}\end{blocLivre}
\begin{blocLivre}{cor-5-9}\begin{corrige}{9}{Deux symboles}\demarche{Retirer d’abord ce qui est connu.}
Les deux triangles valent ensemble $19-5=14$, donc $\triangle=\res7$. Alors $7+\bullet=12$ donne $\bullet=\res5$. Un même symbole cache toujours le même nombre.
\end{corrige}\end{blocLivre}
\begin{blocLivre}{cor-5-10}\begin{corrige}{10}{Le nombre caché}\demarche{Remonter les opérations une à une.}
Avant d’ajouter 3, le produit valait $31-3=28$. Avant de multiplier par 4, le nombre était $28:4=\res7$. Contrôle : $7\times4+3=31$.
\end{corrige}\end{blocLivre}
\begin{blocLivre}{cor-5-11}\begin{corrige}{11}{Les tee-shirts}\demarche{Séparer ce qui dépend du nombre acheté et ce qui est fixe.}
Pour 4 tee-shirts : $4\times12+5=\res{53}$ euros. Pour $N$ tee-shirts, le prix en euros s’écrit $\res{(N\times12)+5}$ : la lettre $N$ représente n’importe quel nombre de tee-shirts.
\end{corrige}\end{blocLivre}
\begin{blocLivre}{cor-5-12}\begin{corrige}{12}{Ciseaux et stylos}\demarche{Retirer la dépense connue avant de partager.}
Les stylos coûtent $3\times2=6$ euros. Les deux paires coûtent $20-6=14$ euros, donc une paire coûte $14:2=\res7$ euros. Contrôle : $2\times7+3\times2=20$.
\end{corrige}\end{blocLivre}
\begin{blocLivre}{cor-5-13}\begin{corrige}{13}{Des carrés d’allumettes}\demarche{Compter le premier motif puis les ajouts identiques.}
Pour 5 carrés : $4+4\times3=\res{16}$ allumettes. Pour 10 carrés : $4+9\times3=\res{31}$. Après le premier carré, chaque nouveau carré partage une allumette avec le précédent.
\label{fig-cor-5-13}\figureGeo[corrige]{motifs-13}
\end{corrige}\end{blocLivre}
\begin{blocLivre}{cor-5-14}\begin{corrige}{14}{Le motif en croix}\demarche{Compter les ajouts depuis la première étape.}
Étape 4 : $5+3\times4=\res{17}$ carreaux. Pour 29 carreaux : $29-5=24$, soit $24:4=6$ ajouts après l’étape 1 : c’est l’étape $\res7$. Contrôle : $5+6\times4=29$.
\label{fig-cor-5-14}\figureGeo[corrige]{motifs-14}
\end{corrige}\end{blocLivre}
\begin{blocLivre}{cor-5-15}\begin{corrige}{15}{Inventer un programme}\demarche{Tester chaque programme sur le même nombre.}
Une solution : multiplier par 5 puis ajouter 1, car $5\times5+1=\res{26}$. Les deux programmes diffèrent : $5\times5+1=26$ mais $(5+1)\times5=\res{30}$. L’ordre des étapes change le résultat.
\end{corrige}\end{blocLivre}
